% Proposed edition Chapter 280. Authored navigation to existing reader entries.
\chapter{Finder by Reasoning Task}
\label{sec:p14-reasoning-task-finder}
\index{reasoning task!finder}\index{method selection}

Start with the question to be answered and the information actually supplied.
A proof, a satisfying assignment, a posterior distribution, an identified
causal effect and an optimal policy answer different questions. Each lookup
below states an input, a useful destination and the conditions governing its
output. Several destinations can be needed for one inquiry: formalize an
observation, compare explanations, then choose a measurement.

The destination gives the full definitions, source presentation, result status
and proof or citation. A link here preserves that standing. A method requiring
a finite input, a supplied graph, a sampling model or a particular calculus
does not acquire a wider guarantee by appearing in this finder.

\begin{longtable}{@{}p{.42\textwidth}p{.50\textwidth}@{}}
\caption{Finder by Reasoning Task}\\
\toprule
Question to start with & First destination in this chapter \\
\midrule
\endfirsthead
\caption[]{Finder by Reasoning Task (continued)}\\
\toprule
Question to start with & First destination in this chapter \\
\midrule
\endhead
\midrule\multicolumn{2}{r}{\emph{continued on next page}}\\
\endfoot
\endlastfoot
What exactly does the assertion or argument say? &
\hyperref[sec:p14-finder-formalizing]{Formalizing an assertion or argument}. \\
Must the conclusion follow, or can it fail? &
\hyperref[sec:p14-finder-proving]{Proving a consequence and finding a counterexample}. \\
Can this problem always be answered effectively, and can this language express it? &
\hyperref[sec:p14-finder-deciding]{Deciding satisfiability, equivalence or expressibility}. \\
What do observations support about an unknown quantity, future event or explanation? &
\hyperref[sec:p14-finder-evidence]{Estimating, predicting, explaining and identifying causes}. \\
What should be chosen, planned, accepted collectively or challenged? &
\hyperref[sec:p14-finder-choosing]{Choosing, planning, aggregating judgments and evaluating arguments}. \\
\bottomrule
\end{longtable}

\section{Formalizing an assertion or argument}
\label{sec:p14-finder-formalizing}

\subsection{Fix the objects, quantifiers and scope}
\label{sec:p14-finder-language}
\textit{Input and question.} An assertion such as \emph{every request has a
response}: what counts as a request, response and response relation, and does
one response serve all requests or depend on the request?
\textit{Use and output.} Specify the signature, domain, interpretation,
quantifier order and free variables; obtain a formula together with its
intended reading. Start with
\hyperref[sec:p01-symbol-registry]{binding and notation},
\hyperref[sec:prop-syntax-semantics-01]{propositional syntax} and
\hyperref[sec:fol-semantics-02]{first-order satisfaction}.
\textit{Conditions.} Ordinary first-order quantifiers range over the declared
nonempty domain. For different kinds of objects, use the explicit
\hyperref[sec:p02-many-sorted]{strict many-sorted presentation}, including its
nonempty disjoint sorts. Translation into a language does not establish that
the chosen interpretation captures the original assertion.

\subsection{Recover an argument from a passage}
\label{sec:p14-finder-reconstruction}
\textit{Input and question.} A passage, its context and supporting records:
what conclusion is offered, by which premises, at what strength?
\textit{Use and output.}
\hyperref[def:p18-scheme-reconstruction-method]{Argument reconstruction}
separates explicit premises, supported attributions of implicit commitments,
and needed but unestablished premises. It returns a reconstruction and a
qualified assessment.
\textit{Conditions.} Preserve qualifications, time, population and the
distinction between describing a fact and recommending an action. Compare
plausible readings; adding an unsupported premise that guarantees the
conclusion does not establish the speaker's argument.
The selected scheme presentations and their distinct critical questions are
at \cref{sec:p18-argument-schemes-critical-questions}
\cite[pp.~271--272, 277--283]{GoddenWalton2007SchemesQuestions}.

\subsection{Represent time, possibility or information}
\label{sec:p14-finder-modality}
\textit{Input and question.} States, possible transitions or information
alternatives, with a claim such as \emph{a response remains possible at every
reachable state}.
\textit{Use and output.} Choose the semantic object before the operators:
\hyperref[def:p08-ctl-semantics]{CTL} describes branching state properties;
\hyperref[def:p08-ctl-star-semantics]{CTL*} keeps state and path formulas
distinct. For a truthful public message, the
\hyperref[def:p08-public-announcement]{public-announcement update} specifies
which worlds remain.
\textit{Conditions.} The temporal entries use their declared total transition
structures; an announcement has its stated truth condition. A modal frame or
information model is supplied mathematical structure. Its accessibility
relation does not establish a causal mechanism or a numerical probability.
Use \cref{sec:p14-modal-correspondence} to compare frame assumptions.

\subsection{Represent exceptions, contradiction or belief change}
\label{sec:p14-finder-defeasibility}
\textit{Input and question.} Are some records inconsistent, are rules intended
to hold normally, or is new information to revise previous commitments?
\textit{Use and output.} These lead respectively to
\hyperref[def:p08-paraconsistency]{nonexplosive consequence},
\hyperref[def:p08-nonmonotonic]{nonmonotonic consequence} and
\hyperref[def:p08-agm]{revision or contraction of belief sets}.
\textit{Conditions.} Select the actual semantics or postulates, rather than a
family name alone. Failure of explosion permits some contradictions without
every conclusion following; it does not certify the records as true. Default
acceptance may change after additional premises. AGM revision concerns its
specified deductively closed belief sets, rather than automatically defining
a probability update. The separate presentations are in
\cref{sec:p08-paraconsistent-and-nonmonotonic}.

\subsection{Represent a question or a dependence between answers}
\label{sec:p14-finder-question-language}
\textit{Input and question.} An explicit list of possible answers, or a claim
that resolving one question resolves another.
\textit{Use and output.}
\hyperref[sec:p08-inquisitive-propositional]{Inquisitive propositional logic}
represents support and resolution by information states.
\hyperref[sec:p18-question-evocation]{PMCE question evocation} instead asks
whether declarative premises cover an answer list while leaving each answer
unsettled.
\textit{Conditions.} Keep these relations distinct. Support of a question
settles it; evocation leaves its individual answers open. The current PMCE
input uses finite propositional premises and an ordered list of at least two
syntactically distinct answers. Neither formalism assigns an experiment's cost
or promises that an answer can be obtained
\cite{Ciardelli2023InquisitivePropositional,Wisniewski2016QuestionEvocation}.

\subsection{Identify the standard intended by rationality}
\label{sec:p14-finder-rationality}
\textit{Input and question.} A claim that an action, belief or person is
rational: does it concern structural coherence, response to actual reasons,
beliefs about reasons, or what motivated the response?
\textit{Use and output.}
\hyperref[def:p11-reasons-diagnostic]{The reasons--rationality diagnostic}
locates the relevant conception and the premise on which the assessment
depends. The separately attributed
\hyperref[def:p11-ri-framework-comparison]{framework comparison}
also distinguishes theoretical from practical, normative from descriptive,
individual from collective, and process from output assessments. For
example, \cref{thm:p11-ri-individual-joint} separates an individually
best response from a jointly beneficial outcome; neither criterion
is inferred from the bare word \emph{rational}.
\textit{Conditions.} The distinctions are interpretive and normative. A
correct calculation under supplied beliefs and utilities does not establish
the truth of those beliefs, the adequacy of the utilities or the authority of
a normative requirement. Read the separately attributed discussion at
\cref{sec:p11-reasons-and-rationality,sec:p11-rationality-intro}.

For the specific structural account, \cref{def:p11-rat105-dimensions}
distinguishes diachronic, intrapersonal and interpersonal assessments.
For an evolutionary functional assessment,
\cref{def:p11-rat106-task-analysis} requires an independently supported
adaptive problem and discriminating mechanism predictions. A cooperation
preference can have the representation of \cref{thm:p11-rat105-value}
without proving that it satisfies every structural requirement. Likewise,
selection for a reproductive function does not establish an epistemic
entitlement or moral obligation \cite{NidaRumelinGutwaldZuber2021,
CosmidesTooby2021Adaptationism}.

\subsection{Use a diagram as a logical representation}
\label{sec:p14-finder-logical-diagrams}
\textit{Input and question.} A diagram, intended set or spatial
interpretation and proposed transformation: what is expressed, and does
the conclusion follow from the represented information?
\textit{Use and output.} \Cref{sec:p12-rationality-13-1} distinguishes
Euler, Venn-I/II and spider representations. Its exact spider reading is
\cref{thm:p12-r131-cardinality}; the licensed fragment has its soundness
proof at \cref{thm:p12-r131-rules}. The checkerboard invariant and
uniform ell decomposition separately address arithmetic representations.
\textit{Conditions.} Missing zones, shading, linked alternative witness
locations and distinct spiders contribute different information. Incidental
shape or an attractive drawing does not establish an inclusion or a proof.
A valid worked transformation is not full-system completeness or evidence
that every person reasons better with a picture \cite{Jamnik2021LogicalDiagrams}.

\subsection{Answer a spatial or durative temporal query}
\label{sec:p14-finder-relational-models}
\textit{Input and question.} Interpreted spatial or interval premises and
a proposed conclusion: is it possible in some satisfying configuration,
necessary in every configuration, or inconsistent with the premises?
\textit{Use and output.} \Cref{def:p19-rat132-task} separates these tasks;
\cref{thm:p19-rat132-endpoints} gives the finite endpoint-order search for
proper intervals. The car-order example at \cref{thm:p19-rat132-cars}
exhibits all three satisfying orders rather than only the preferred one.
\textit{Conditions.} Fix the reference frame, relation interpretation and
duration assumptions. A sound local propagation step is not global
consistency: \cref{thm:p19-rat132-propagation} retains that boundary.
An interval representation does not itself identify a person's cognitive
process, and temporal priority does not imply causation
\cite{Ragni2021SpatialTemporal,Allen1983TemporalIntervals}.

\subsection{Locate a historical name or source presentation}
\label{sec:p14-finder-history}
\textit{Input and question.} A name, date, theorem attribution or unfamiliar
historical terminology: which work and technical entry should be consulted?
\textit{Use and output.}
\hyperref[sec:p13-chronology]{The selected chronology} and
\hyperref[sec:p13-biographies]{biographical capsules} give publication and
contributor leads. Follow the cited work and then the relevant technical
entry for its actual language, hypotheses and result.
\textit{Conditions.} These are selected reference records. A publication date
does not by itself establish priority for every related formulation, and a
shared theorem name does not identify editions, proof variants or semantic
conventions.

\section{Proving a consequence and finding a counterexample}
\label{sec:p14-finder-proving}

\subsection{Test a finite classical propositional argument}
\label{sec:p14-finder-propositional}
\textit{Input and question.} Finitely many propositional premises and a
proposed conclusion: is there a valuation making all premises true and the
conclusion false?
\textit{Use and output.} Search the conjunction of the premises and negated
conclusion. A satisfying assignment is a counterexample; unsatisfiability
establishes consequence under the classical truth conditions. Use
\hyperref[def:p02-cpl-tableau]{analytic tableaux} or convert to clauses and
use \hyperref[def:p12-dpll]{the displayed DPLL procedure}.
\textit{Conditions.} The DPLL termination and exact-output guarantee is
\cref{thm:p12-dpll-correctness}, for its finite clauses, unit propagation and
binary branching. Signature-expanding transformations require their own
model-reconstruction relation. A returned solver word alone is not the
\hyperref[def:p12-sat-certificate-boundary]{model or refutation certificate}.

\subsection{Find a common instance of finite terms}
\label{sec:p14-finder-syntactic-unification}
\textit{Input and question.} A finite multiset of equations between finite
first-order terms: is there a simultaneous substitution making both sides
literally identical, and what is its most general answer?
\textit{Use and output.} The retained-equation sheet and its extraction
are \cref{def:p12-ue-sheet-run,thm:p12-ue-extraction}. The distinct
removing accumulator and its composition invariant are
\cref{thm:p12-ue-accumulator}. Both either return an idempotent most
general unifier or certify failure under their finite-term assumptions.
\textit{Conditions.} Apply an enabled rule and continue to genuine
irreducibility, not merely an equation on which a rule is disabled. Keep
the simultaneous substitution and composition direction fixed.
The occurs check rejects \(x\doteq f(x)\) as a finite-term equation;
removing it changes the term domain. Most general answers are unique
up to the variable renaming of \cref{thm:p12-ue-renaming}; factor
substitutions need not be unique. Equational and higher-order unification
are separate problems, not guarantees inherited from this procedure
\cite{MartelliMontanari1982,BaaderSnyder2001}.

\subsection{Solve equations modulo a declared theory}
\label{sec:p14-finder-equational-unification}
\textit{Input and question.} Finite terms, an equational theory and a
problem-variable set: which substitutions solve the equations modulo the
theory, and does the returned family cover all solutions?
\textit{Use and output.} \Cref{def:p12-eu-generality} fixes comparison
only on the problem variables. The separately specified G, B and N
calculi and their completeness hypotheses are in
\cref{sec:p12-n-c-035-05}; the commutative constructions distinguish
ACU, ACUI and Abelian groups.
\textit{Conditions.} There can be multiple incomparable members of a
minimal complete set, infinitely many such members, or no minimal complete
set at all. A syntactic MGU does not solve a nonempty equational theory.
The ordered closure, unrestricted closure and strategy-restricted systems
have different hypotheses and are not one effective procedure
\cite{BaaderSnyder2001,DershowitzPlaisted2001Rewriting}.

\subsection{Check a proof term in its exact type system}
\label{sec:p14-finder-type-checking}
\textit{Input and question.} A context, annotated term, alleged type and
declared conversion relation: does the term establish that typing judgment?
\textit{Use and output.} \Cref{def:p12-har-three-problems} distinguishes
checking, synthesis and inhabitation. The exact beta-only source
lambdaHOL checker at \cref{thm:p12-har-checker-conditional} uses the
full normalization and transport proofs of
\cref{thm:p12-har-cc-normalization,thm:p12-har-hol-normalization}.
\textit{Conditions.} Decidable checking does not imply decidable search
for a proof. Neither the source's name nor a modern assistant's name
extends this theorem to delta/iota reductions, inductive additions,
strong sums or different universe rules. An implementation still needs
its own correctness and trust account
\cite{BarendregtGeuvers2001,Geuvers1994Report9450}.

\subsection{Refute a first-order conjecture}
\label{sec:p14-finder-first-order}
\textit{Input and question.} A finite first-order background and a closed
conjecture: can their counterexample conditions be refuted?
\textit{Use and output.}
\hyperref[sec:p12-small-clause-standard]{Standard clausification} records
the passage to clauses; \hyperref[def:p12-resolution-search]{resolution search}
records the parent clauses, substitutions and rule instances of a refutation.
\textit{Conditions.} Skolemization preserves satisfiability using a fresh
expanded signature and the stated expansion/reduct assumptions; it is not
same-signature equivalence. The resolution completeness entry
\cref{thm:fol-resolution} concerns its equality-free finite-clause calculus
with factoring. A search procedure must meet its own fairness and redundancy
conditions to find every available refutation. An unsuccessful
bounded search proves neither satisfiability nor failure of consequence.
Equality and ordered refinements require their own registered rules.

\subsection{Check a proof or remove its detours}
\label{sec:p14-finder-proof-transformation}
\textit{Input and question.} A finite derivation: is every inference licensed,
and can the derivation be transformed while preserving its conclusion?
\textit{Use and output.} Consult \cref{sec:p14-rule-sheets} for the exact
calculus. For full first-order intuitionistic natural deduction with empty
discharge classes, use
\cref{thm:p07-nj-vacuous-normalization,sheet:p07-nj-vacuous}; the result
supplies a terminating normalization strategy.
\textit{Conditions.} Preserve open assumptions, binding, discharge and
eigenvariable restrictions. The separately registered nonvacuous
\hyperref[sheet:nj]{NJ presentation} has the obstruction in
\cref{thm:p07-nj-nonvacuous-obstruction}, despite the consequence bridge in
\cref{thm:p07-nj-vacuous-consequence-bridge}. For arrow-only typed terms,
\cref{thm:p07-stlc-strong-normalization} supplies its own stronger
all-reduction-sequences claim. These are different exact targets.

\subsection{Prove identities or compare terms by normalization}
\label{sec:p14-finder-equations}
\textit{Input and question.} Equations assumed as identities of every algebra
under consideration, and a target equation.
\textit{Use and output.}
\hyperref[sec:p10-equational-logic]{Equational deduction} uses substitution,
congruence and equality rules to derive an identity. A successful
\hyperref[def:p12-completion]{completion construction} can supply a convergent
rewrite presentation, after which normal forms compare terms.
\textit{Conditions.} Identities quantify over every valuation in each
nonempty algebra; a local equation true at one valuation has a different
scope. The normalization criterion
\cref{thm:p12-completion-correctness} requires a finite completed run,
no remaining equations, joinable critical pairs and rules decreasing under
the fixed reduction ordering. An orientation failure or an infinite run does
not refute the proposed identity.

\subsection{Verify a program or a finite transition model}
\label{sec:p14-finder-verification}
\textit{Input and question.} Either a command with precondition and
postcondition, or an explicit finite transition structure with a temporal
formula.
\textit{Use and output.}
\hyperref[def:p09-hoare-triple]{Hoare triples} distinguish partial correctness
from total correctness. \hyperref[def:p09-ctl-labelling]{CTL labelling}
returns exactly the satisfying states of the supplied finite total structure
by \cref{thm:p09-ctl-labelling}.
\textit{Conditions.} Partial correctness concerns terminating executions;
termination requires the total-correctness obligations. The Hoare entry fixes
a deterministic while language and its assertion semantics. CTL verification
depends on the represented states, transitions and labels. Neither result
alone establishes that a model includes every behavior of a physical or
software implementation.

\subsection{Find a least-cost path}
\label{sec:p14-finder-path}
\textit{Input and question.} A state graph, start, goals, nonnegative edge
costs and heuristic: which reachable goal has the cheapest path?
\textit{Use and output.}
\hyperref[def:p12-hsearch-algorithm]{A* with duplicate detection and reopening}
returns a path. On a finite graph it terminates; an admissible heuristic
gives least cost under \cref{thm:p12-hsearch-finite,thm:p12-hsearch-optimal}.
\textit{Conditions.} Return on goal selection, retain valid path witnesses
and reopen improved states as specified. Consistency licenses the separate
no-reopening result. On infinite graphs use the extra assumptions of
\cref{thm:p12-hsearch-infinite}; even that theorem does not promise a negative
answer on an infinite unsolvable graph
\cite{DechterPearl1985BestFirst}.

\subsection{Transform a difficult problem or work backwards}
\label{sec:p14-finder-backward}
\textit{Input and question.} Unknowns, data, domain, conditions and a desired
output: can a simpler problem supply a route to a solution?
\textit{Use and output.}
\hyperref[sec:p12-heuristic-transformations-and-backward-reasoning]{Heuristic
transformations and backward reasoning} separate necessary, sufficient and
equivalent transformations, with reconstruction of solutions.
\textit{Conditions.} A backward chain must justify the forward implications
that eventually establish the target. Squaring, division, specialization and
relaxation retain their domain and reconstruction conditions. A candidate
surviving a necessary condition still needs verification in the original
problem. The worked methods give search guidance; they provide no universal
discovery or termination theorem
\cite{Polya2004HowToSolveItSelected}.

\section{Deciding satisfiability, equivalence or expressibility}
\label{sec:p14-finder-deciding}

\subsection{Ask for the exact decision problem}
\label{sec:p14-finder-decision-input}
\textit{Input and question.} A language and class of inputs: is the requested
property decidable, semidecidable or subject to a known complexity bound?
\textit{Use and output.} Fix whether the question is validity, satisfiability,
consequence, equivalence or model checking; fix the structures, encoding and
which objects are inputs. Then use
\hyperref[sec:p14-decision-table]{the decision-problem table} and its linked
statements. \hyperref[sec:decision-hierarchy-02]{The arithmetical hierarchy}
provides finer classifications for effectively coded sets.
\textit{Conditions.} General first-order validity is enumerable but
undecidable by \cref{thm:fol-validity-undecidable,thm:fol-validity-enumerable}.
This does not classify each restricted fragment. The current precise
fragment boundaries are in \cref{tab:p02-decidable-fragments}.

\subsection{Distinguish a finite model from arbitrary models}
\label{sec:p14-finder-finite-models}
\textit{Input and question.} A first-order sentence: must it hold in all
structures, in all finite structures, or only in one supplied structure?
\textit{Use and output.}
\hyperref[def:p09-finite-validity]{Finite satisfiability and validity}
fix the first two finite-model problems; an explicit finite structure instead
allows evaluation by its satisfaction clauses.
\textit{Conditions.} Exhausting structures up to one chosen size proves only
that bounded absence. For every fixed finite relational signature containing
a binary relation, \cref{thm:trakhtenbrot} gives enumerable but undecidable
finite satisfiability and nonenumerable finite validity
\cite{Trakhtenbrot1950}. A completeness theorem for unrestricted first-order
semantics cannot be transferred to finite validity.

\subsection{Compare countable structures up to isomorphism}
\label{sec:p14-finder-countable-isomorphism}
\textit{Input and question.} Two nonempty at most countable structures in
the same countable language: do they have the same isomorphism type, and
what does a count of models actually count?
\textit{Use and output.} The coding and infinitary construction at
\cref{sec:p03-n-a-004-02} distinguishes labeled copies from isomorphism
classes. \Cref{thm:p03-morley-scott-sentence} constructs a sentence
characterizing one structure among at most countable structures;
\cref{thm:p03-morley-approximant-isomorphism} characterizes isomorphism
by all countable back-and-forth levels.
\textit{Conditions.} The construction is not an effective finite decision
procedure for arbitrary presentations. An individual isomorphism class
being Borel does not make the global isomorphism relation Borel. Nor does
a perfect family for an incomplete theory decide the spectrum conjecture
for complete theories; the exact distinctions are in
\cref{thm:p03-morley-isomorphism-complexity,thm:p03-morley-perfect-example}.

\subsection{Use dimension to classify strongly minimal models}
\label{sec:p14-finder-strongly-minimal-dimension}
\textit{Input and question.} A specified strongly minimal definable set
with its parameter base: what does algebraic-closure dimension determine?
\textit{Use and output.} \Cref{thm:p03-bl-pregeometry,thm:p03-bl-dimension}
give exchange, bases and their cardinal invariant. Under the additional
whole-universe strongly minimal hypothesis,
\cref{thm:p03-bl-whole-models} classifies models by basis cardinality.
The fixed-characteristic field illustration is
\cref{thm:p03-bl-acf-geometry}.
\textit{Conditions.} A strongly minimal proper subset need not control
the rest of its model: \cref{def:p03-bl-extra-sort} separates equal subset
dimensions from unequal complement cardinalities. General categorical
coordinatization requires the further theorem, not the definition of a
basis. Rank, degree and basis dimension are different invariants
\cite{Morley1965,BaldwinLachlan1971}.

\subsection{Compare formulas or solve constraints in a theory}
\label{sec:p14-finder-equivalence-smt}
\textit{Input and question.} Two Boolean formulas, or constraints involving
interpreted operations such as arithmetic and equality.
\textit{Use and output.} For Boolean equivalence, compare truth functions;
\hyperref[sec:p12-sat-and-certificates]{the SAT and decision-diagram chapter}
gives finite search and ordered decision-diagram representations. For
interpreted constraints, \hyperref[def:p12-smt-problem]{an SMT instance}
declares the background theory, sorts and input fragment.
\textit{Conditions.} Boolean proxies forget the interpretation of theory
atoms, so a satisfying proxy assignment still needs theory consistency.
Canonical reduced ordered diagrams use the same variable order.
\cref{tab:p12-smt-theories,thm:p12-nelson-oppen} retain the conditions
on theory decision procedures and combination. No guarantee applies to
arbitrary mixed theories merely because an interface accepts their syntax.

\subsection{Show that a language cannot express a property}
\label{sec:p14-finder-expressibility}
\textit{Input and question.} A class of structures or behavior property and
a proposed defining language.
\textit{Use and output.} For first-order definability on finite structures,
\hyperref[sec:p03-ef-games]{Ehrenfeucht--Fra\"iss\'e games} convert a
quantifier-rank bound into a comparison of structures. The criterion
\cref{thm:p03-ef-inexpressibility} uses an inside/outside pair for every
finite rank.
\textit{Conditions.} One fixed finite-round win gives only bounded-rank
agreement. The displayed result fixes a finite relational signature; functions
need a separate treatment. For modal or temporal comparisons consult
\cref{sec:p14-expressive-table,thm:p09-ltl-ctl-boundary}, retaining state/path
typing and the model class
\cite{Ehrenfeucht1961,Libkin2004,EmersonHalpern1986}.

\subsection{Locate an independence or strength claim}
\label{sec:p14-finder-strength}
\textit{Input and question.} A statement, formal theory and proposed
comparison: is the statement unprovable there, or is one theory stronger in
a specified sense?
\textit{Use and output.}
\hyperref[sec:p05-independence-catalog]{The independence catalog} separates
standard-model truth, formalization and unprovability.
\hyperref[sec:p06-consistency-strength]{The consistency-strength map}
separates inclusion, interpretation, conservation and relative consistency.
\textit{Conditions.} Keep the exact axioms, proof predicate, base theory and
consistency or soundness assumptions. A failed proof search is not an
independence proof. An inclusion of theories is not automatically a strict
increase in consistency strength; an ordinal analysis is another comparison.
Use \cref{sec:p14-strength-table} for the corresponding retrieval table.

\section{Estimating, predicting, explaining and identifying causes}
\label{sec:p14-finder-evidence}

\subsection{Update uncertainty after an observation}
\label{sec:p14-finder-update}
\textit{Input and question.} A prior, sampling model and observation: how
does the specified joint model condition the unknown quantity?
\textit{Use and output.}
\hyperref[thm:p15-finite-bayes]{Finite Bayes conditioning} handles a finite
partition; \cref{thm:p15-bayes-posterior-kernel} constructs the posterior kernel
for the displayed dominated experiment.
\textit{Conditions.} The prior and likelihood are inputs; the observation
must meet the positive, finite normalizing condition for the displayed
formula. Continuous conditioning requires a specified density or kernel
version. Successive data use their conditional likelihoods: a copied report
is not a second independent measurement. The exact contrast is
\cref{thm:p15-bayes-flag-calculation}.

\subsection{Assess a credence norm or select an update rule}
\label{sec:p14-finder-credence-dynamics}
\textit{Input and question.} An event algebra, interpreted credences and
a description of what is learned: which coherence claim, update construction
or proposed rational requirement is being assessed?
\textit{Use and output.} \Cref{sec:p15-rationality-4-1} separates the
probability calculus from probabilism and its betting, representation and
accuracy justifications. For new weights on a finite partition, use
\cref{def:p15-rat41-jeffrey,thm:p15-rat41-jeffrey}; for a constrained
minimal-change proposal, specify the divergence and its orientation as in
\cref{thm:p15-rat41-minimal-update}.
\textit{Conditions.} A representation or accuracy calculation does not
prove the normative authority of its premises. Jeffrey updates require
positive prior cell masses and retained within-cell laws; sequential resets
need not commute or preserve every earlier marginal constraint. Variation
minimizers need not be unique even where KL and Hellinger minimizers are.
Approximate coherence depends on the stated score, distance and stake
normalization in \cref{thm:p15-rat41-graded}
\cite{HajekStaffel2021Dynamics}.

\subsection{Find the coherent range of an uncertain conclusion}
\label{sec:p14-finder-probability-logic}
\textit{Input and question.} Interpreted propositions, possible cases and
conditional assessments or intervals: which conclusion values are coherently
extendible from those premises?
\textit{Use and output.} \Cref{sec:p15-rationality-4-4} fixes primitive
conditional events and the finite coherent-extension question. Its recursive
conditioning strata and sharp propagation bounds preserve null-but-possible
antecedents. The finite ratio contrasts in
\cref{sec:p19-rationality-4-6,sec:p19-rationality-6-2} concern separately
declared interpretations and positive-denominator conditions.
\textit{Conditions.} A probability-zero event need not be impossible, and
a primitive conditional assessment need not be a defined ratio. A feasible
range is not a prediction of a participant's particular report. Fixed
numerical constraints and revision of a default's conditioning context are
different operations. Nested conditional random quantities need their own
semantics; the nonnested construction is not their proof
\cite{Pfeifer2021ProbabilityLogic,OverCruz2021Suppositional}.

\subsection{Choose a representation of acceptance or imprecise belief}
\label{sec:p14-finder-belief-representations}
\textit{Input and question.} Accepted propositions, graded confidence,
information sets or buying/selling assessments: which objects and combination
or updating question are actually supplied?
\textit{Use and output.} \Cref{sec:p15-rationality-4-7} separates necessity
and possibility, negative ranks, focal-set belief functions, pignistic
projection and lower/upper previsions. Dynamic ranking operations are
developed separately in \cref{sec:p15-rationality-5-3}.
\textit{Conditions.} A static transform does not transfer every update law.
Independent-source combination requires positive nonconflict probability;
pignistic projection need not recover the original focal masses. The binary
minimum rank axiom alone does not ensure a nonempty intersection of every
believed event on an arbitrary infinite space. Two-cell ranking updates
require the stated finite cell ranks and specify a target firmness, not an
unqualified increment \cite{DuboisPrade2021RepresentingBelief}.

\subsection{Separate evidential relevance from posterior confidence}
\label{sec:p14-finder-evidential-relevance}
\textit{Input and question.} A prior joint table, target hypothesis, evidence
and chosen measure: does the evidence support the target, and how is that
support compared with another item?
\textit{Use and output.} \Cref{sec:p15-rationality-4-3} gives signed support,
screening, sufficiency and measure comparisons at their explicit domains.
\Cref{sec:p16-rationality-4-2} adds reliability and search models with their
actual conditional-independence assumptions.
\textit{Conditions.} A high posterior can accompany negative support, and
a lower conjunction posterior can accompany greater relevance. Fixed-target
rankings do not license unrestricted source/target reversal. Reports sharing
reliability need not supply independent likelihood factors; failure to find
an alternative differs from its nonexistence. Numerical parameters require
justification rather than being supplied by a diagram
\cite{Merin2021Relevance,Hartmann2021BayesNets}.

\subsection{Predict a future observation}
\label{sec:p14-finder-prediction}
\textit{Input and question.} A posterior and a specified law for a future
observation: what predictive distribution follows?
\textit{Use and output.}
\hyperref[sec:p15-bayes-prediction]{Posterior prediction} integrates the future
kernel over the posterior, retaining uncertainty about its parameter.
\textit{Conditions.} Declare whether the future is conditionally independent
of past data given that parameter, and whether the observation mechanism
changes. This is a probability calculation within the supplied model, not a
proof that future data will follow it. For assessment on new observations,
also consult \cref{sec:p15-bounded-concentration-learning,sec:p19-evidence-methods}.

\subsection{Test a hypothesis or estimate with coverage}
\label{sec:p14-finder-testing}
\textit{Input and question.} A sampling experiment, null and alternative,
error level, and target parameter: what test or confidence set has the
specified repeated-sampling guarantee?
\textit{Use and output.}
\cref{thm:p15-neyman-pearson} gives the most-powerful comparison of two
specified dominated laws.
\cref{thm:p15-confidence-test-inversion} turns level-controlled tests into
confidence sets; the binomial construction is
\cref{thm:p15-binomial-confidence-coverage}.
\textit{Conditions.} Preserve the sampling and randomization scheme, nuisance
parameter quantifiers and measurability. Nonrejection does not prove the null.
Coverage of a random procedure is different from posterior probability for
the parameter after observing its realized set
\cite{NeymanPearson1933Tests,Neyman1937Confidence}.

\subsection{Bound estimation error or select a predictor}
\label{sec:p14-finder-learning}
\textit{Input and question.} Bounded observations and a specified sampling
scheme, possibly with a finite candidate class: how far can observed averages
deviate from expectations?
\textit{Use and output.}
\hyperref[sec:p15-bounded-concentration-learning]{Bounded concentration and
finite-class learning} provide probability bounds, uniform comparisons and
validation constructions.
\textit{Conditions.} Use the actual range, independence or separately proved
dependence assumptions. A bound for one fixed candidate is not automatically
a bound for the candidate selected after examining the same data; use the
uniform or independent-validation result appropriate to that selection.
Concentration controls sampling deviation under its law, not bias caused by
an unrepresentative population or an incorrect response model
\cite{Hoeffding1963BoundedSums}.

\subsection{Ask what a fitted model fails to reproduce}
\label{sec:p14-finder-model-criticism}
\textit{Input and question.} Observed data, posterior, replication scheme and
discrepancy: which features look unusual relative to the fitted model?
\textit{Use and output.}
\hyperref[def:p15-ppc-simulation]{Posterior predictive comparison} pairs
replicated and observed discrepancies under the declared joint law;
\cref{thm:p15-ppc-joint} identifies the resulting reference probability.
\textit{Conditions.} Fix the replicated experiment and any conditioned
auxiliary information. Parameter-dependent discrepancies use the same
posterior parameter draw in both terms. The calibration examples in
\cref{sec:p15-ppc-calibration} show why this probability is not automatically
a uniformly calibrated frequentist \(p\)-value. Agreement on one discrepancy
does not validate every feature of a model.

\subsection{Calculate a stipulated intervention}
\label{sec:p14-finder-intervention}
\textit{Input and question.} Structural equations, disturbance distribution
and variables to be externally set: what follows after replacing those
equations?
\textit{Use and output.}
\hyperref[def:p16-causal-finite-calculation]{Finite intervention calculation}
returns the interventional distribution of the specified acyclic structural
model.
\textit{Conditions.} Equation replacement is part of the causal model's
definition; ordinary conditioning is a different operation. Product and
truncated-factorization formulas need the stated disturbance independence.
\cref{thm:p16-causal-confounded-example,thm:p16-causal-nonidentification}
give concrete association/intervention and observational-equivalence
boundaries. A calculation from supplied structural equations does not by itself
show that those equations can be recovered from observations.

\subsection{Identify a causal effect from an observational law}
\label{sec:p14-finder-identification}
\textit{Input and question.} A causal graph or model class, observed law,
treatment and response: do all compatible models agree on the requested effect?
\textit{Use and output.}
\cref{thm:p16-finite-back-door-adjustment} supplies a sufficient adjustment
formula. For the finite positive-law latent-input setting,
\hyperref[def:p16-id-algorithm]{the ID algorithm} returns an expression or a
structural failure record, with termination and successful-return soundness in
\cref{thm:p16-id-termination,thm:p16-id-success-soundness}.
\textit{Conditions.} Check the actual graph, observed variables, independence
and overlap assumptions. A graph is not inferred merely by supplying a
positive table. Zero support is outside that execution theorem.
Identification from a law and estimation of that law from a sample are
separate obligations; an identification formula alone gives no sampling
error bound.

\subsection{Find assumptions that would explain an observation}
\label{sec:p14-finder-abduction}
\textit{Input and question.} A finite propositional background, permitted
positive hypothesis atoms and observed atoms: which allowed assumptions are
consistent with the background and sufficient for the observations?
\textit{Use and output.}
\hyperref[def:p16-abduction-enumeration]{Finite logical abduction}
enumerates explanations and its subset-minimal and minimum-cardinality
families; \cref{thm:p16-abduction-enumeration} gives termination and exact
output.
\textit{Conditions.} Perform both consistency and entailment tests. The
permitted hypothesis set is part of the input, and subset minimality differs
from minimum size. Explanatory sufficiency does not show that the added
assumptions are true or uniquely responsible:
\cref{thm:p16-abduction-truth-boundary,def:p16-abduction-causal-boundary}.

\subsection{Diagnose an inconsistency in a modeled system}
\label{sec:p14-finder-diagnosis}
\textit{Input and question.} A system description, observations and finitely
many component-normality assumptions: which minimal sets of those assumptions
must be withdrawn to restore consistency?
\textit{Use and output.}
\hyperref[sec:p16-consistency-based-diagnosis]{Consistency-based diagnosis}
uses conflicts and hitting sets to characterize the diagnoses.
\textit{Conditions.} The search result with implicit conflict discovery
requires the consistency access of \cref{thm:p16-diag-oracle}; finitely many
component names do not decide an arbitrary first-order background.
Consistency with observations need not entail them. Fault probability and
specific faulty behavior require additional premises, as stated in
\cref{thm:p16-diag-abduction-boundaries,thm:p16-diag-probability-boundary}.

\subsection{Use a numerical result or simulation as evidence}
\label{sec:p14-finder-numerical}
\textit{Input and question.} A computed value, target quantity, input
uncertainty and justified error bounds: which sign, ordering or tolerance
claim is supported?
\textit{Use and output.}
\cref{thm:p12-numerical-certification} converts an enclosing interval into
conditional sign and order certificates.
\hyperref[sec:p12-simulation-results-as-evidence]{Simulation evidence}
connects numerical verification, reference comparison and intended use.
\textit{Conditions.} Include input and model error as well as arithmetic
error. An interval containing zero does not certify a strict sign. Agreement
with an exact mathematical reference does not establish agreement with a
physical target; finite comparisons alone do not supply an unrestricted
prediction bound, as shown in \cref{thm:p12-simulation-finite-comparisons}.

\subsection{Evaluate a judgmental heuristic or diagnose an error}
\label{sec:p14-finder-judgmental-heuristics}
\textit{Input and question.} A task, interpreted information, reported
responses and a proposed search or choice process: is there an error
under the stated benchmark, and does that error identify its mechanism?
\textit{Use and output.} The comparison at
\cref{sec:p19-rationality-2-4} separates coherence, environmental accuracy
and resource costs. \Cref{thm:p19-r24-exposure} shows how selected exposure
can reverse a frequency ordering without biased retrieval;
\cref{thm:p19-r24-priority-invariance} gives a falsifiable restriction of
one precisely specified stopping procedure.
\textit{Conditions.} An outcome shared by different procedures does not
identify one of them. Recognition use is not the same measure as answer
accuracy, and the less-is-more comparison requires the fixed conditional
validities in \cref{thm:p19-r24-recognition-mixture}. Equal probabilities
of exact strings do not force equal observed counts; use the sampling
and conditioning distinctions in
\cref{thm:p19-r24-strings,thm:p19-r24-counts}. An empirical choice pattern
does not by itself establish a normative requirement
\cite{FiedlerPragerMcCaughey2021Heuristics}.

For the source-specific history of why these comparisons were introduced,
\cref{sec:p19-heuristics-program-history} distinguishes the three waves,
their rival process accounts and changes of task interpretation.
\Cref{thm:p19-hph-patterns,thm:p19-hph-support} separate exact sampling
benchmarks from description-sensitive support; the normalized support ratio
does not itself prove a psychological unpacking effect.

\subsection{Compare cognitive explanations of the same response}
\label{sec:p14-finder-cognitive-explanations}
\textit{Input and question.} Task wording, interpreted premises, responses,
timings or neural measurements and explicit rival predictions: what does the
evidence discriminate?
\textit{Use and output.} The distinct accounts at
\cref{sec:p19-rationality-2-5,sec:p19-rationality-3-2,
sec:p19-rationality-3-3,sec:p19-rationality-4-5,
sec:p19-rationality-5-4,sec:p19-rationality-1-4,
sec:p19-rationality-7-2,sec:p19-rationality-7-3} separate process proposals,
formal benchmarks and source-specific evidence records.
\textit{Conditions.} A correct answer, response time, fitted sample parameter
or activation contrast does not identify a unique algorithm. Repeated sessions
of the same participants are not independent populations. A revision choice
does not establish every AGM postulate, and an anatomical contrast does not
by itself distinguish conflict detection from resolution. Specify the rival's
excluded response and the actual design, denominator and uncertainty.
\Cref{sec:p11-ri-process-output} separately identifies normative output,
descriptive output, normative process and descriptive process questions:
agreement on one answer does not establish agreement on resource costs,
the implemented procedure or the authority of the benchmark.

For a Bayesian cognitive account, \cref{def:p19-bci-assessment} requires
the hypothesis representation, observation model, prior and response rule,
not Bayes's identity alone. The finite posterior-sampling convergence
result in \cref{thm:p19-bci-iterated} has its own mixing assumptions;
it does not identify the procedure of every human transmission chain.
For neural evidence, \cref{def:p19-ncm-use} separates measurement,
localization, necessity, perturbation and mechanistic identification.
\Cref{def:p19-ncm-dissociation,thm:p19-ncm-score} keep dissociation
patterns and physiological scores distinct from universal capacity claims.

\subsection{Select the causal judgment being tested}
\label{sec:p14-finder-causal-cognition}
\textit{Input and question.} Learning data, graph instructions, a response
question and observed judgments: is the task prediction, diagnosis,
causal power, structure learning, production attribution or intervention?
\textit{Use and output.} \Cref{def:p19-rat72-task-distinction} separates
the targets; \cref{def:p19-rat73-lookup} selects a diagnostic calculation
and its assumptions. The two wine-cask experiments distinguish learning
within a supplied graph from discovering that graph. Presence and
production have the different quantities in
\cref{thm:p19-rat73-elementary}.
\textit{Conditions.} A common numerical prediction does not establish
the same representation. Training-weighted diagnostic scores need not
be the conditional of a mixture joint law
(\cref{thm:p19-rat73-averaging}). A successful stipulated calculation
does not establish a mental algorithm, clinical performance or transfer
from instructed student tasks to other populations.

\subsection{Predict an action from a valued conditional}
\label{sec:p14-finder-utility-conditionals}
\textit{Input and question.} A conditional connecting an action and a
valued outcome, together with information about the agents: does it support
a fact, an action prediction, an attributed preference or a recommendation?
\textit{Use and output.} \Cref{def:p19-r64-consultation} separates these
tasks and states the additional knowledge, ability, alternatives, costs
and outcome assumptions. \Cref{thm:p19-r64-choice-gap} supplies the
explicit two-act comparison, not a prediction from the conditional alone.
\textit{Conditions.} Values are agent-relative; a recommendation attributed
to another person is not the observer's endorsement. Premise credibility
is different from validity of modus ponens. A diagnostic association is
not an effect of choosing an action, as
\cref{thm:p19-r64-diagnostic-failure} demonstrates. A mean utility rating
is neither an outcome probability nor the same belief held by every
participant \cite{Bonnefon2021Utility}.

\section{Choosing, planning, aggregating judgments and evaluating arguments}
\label{sec:p14-finder-choosing}

\subsection{Separate action assumptions from a social explanation}
\label{sec:p14-finder-social-models}
\textit{Input and question.} Actors, preferences, information, institutions
and a proposed group outcome: which micro--macro links make the prediction?
\textit{Use and output.} \Cref{def:p17-rat104-assumption-types,
def:p17-rat104-coleman} separate rationality, motivation, information,
bridge assumptions and transformation rules. The volunteer game at
\cref{thm:p17-rat104-vod-equilibria} yields different provision under
sole-volunteer and symmetric independent equilibria.
\textit{Conditions.} Transitivity does not identify motives or select an
equilibrium. The cooperative profile is efficient only over its stated
independent-mixture domain, and its group-size effect needs the criterion
of \cref{thm:p17-rat104-provision-correction}. For an empirical economic
comparison, \cref{def:p19-r94-target} separates preference axioms,
representation, auxiliary assumptions and the actual evidence record
\cite{Raub2021RationalChoiceSocialSciences,DhamiAlNowaihi2021Economics}.

\subsection{Choose among lotteries}
\label{sec:p14-finder-choice}
\textit{Input and question.} A finite menu of lotteries over specified
consequences and a preference representation: which available prospect is best?
\textit{Use and output.}
\hyperref[def:p17-expected-utility-choice]{Expected-utility comparison}
returns the maximizing prospects. The finite-prize representation theorem
is \cref{thm:p17-finite-vnm}.
\textit{Conditions.} Its preference assumptions include the stated ordering,
mixture independence and continuity. Utilities of consequences are cardinal
up to the allowed positive affine change; an ordinal ranking of sure outcomes
alone does not fix the ranking of risky prospects. Expected money is the
criterion only when the supplied utility warrants it. The theorem neither
chooses a person's ends nor estimates the probabilities.

\subsection{Compare ambiguity-sensitive decision criteria}
\label{sec:p14-finder-ambiguity-criteria}
\textit{Input and question.} Acts, consequences, utilities and a precise set,
capacity, confidence hierarchy or second-order belief description: which
criterion evaluates the supplied menu?
\textit{Use and output.} \Cref{sec:p17-rationality-8-4} compares maxmin,
alpha-maxmin, unanimity, E-admissibility, Choquet/core, confidence-dependent,
variational, source-weighted and smooth-ambiguity constructions.
\textit{Conditions.} These are not names for one minimum operator. Menus,
act-dependent inputs, utility transformations and conditional-plan assumptions
matter. A representation theorem is conditional on its axioms; a static
preference does not prescribe every later conditional choice. Empirical
source-dependence findings retain their competing representative-agent and
individual-level analyses.

\subsection{Fit or compare a reference-dependent prospect model}
\label{sec:p14-finder-prospect-model}
\textit{Input and question.} Described lotteries, a reference point, value
function and weighting rules: which choices does the specified predictive
model rank above others?
\textit{Use and output.} \Cref{sec:p17-rationality-8-3} separates original
editing/evaluation from cumulative rank-dependent weighting and supplies
reference, curvature, coalescing and sampling contrasts.
\textit{Conditions.} General expected utility does not require concavity.
Mixed cumulative decision weights need not sum to one and are not an ordinary
probability law. The reference and fitted parameters are inputs, not
automatically observed mental processes. A descriptive fit is not a normative
requirement or a guarantee of held-out prediction.

\subsection{Choose a sequence of actions under uncertainty}
\label{sec:p14-finder-sequential}
\textit{Input and question.} States, available actions, transitions, rewards,
information and horizon: which policy maximizes the declared return?
\textit{Use and output.}
\cref{thm:p17-sequential-backward} gives backward recursion and an optimal
deterministic Markov policy for the displayed finite-state, finite-action,
finite-horizon model. The discounted continuation appears in
\cref{sec:p17-sequential-decisions}.
\textit{Conditions.} The transition law must hold under the specified
history and action, and the reward and information timing must match.
For hidden state, use
\cref{sec:p17-partial-observation-decisions,thm:p17-astrom-information-state}
to justify the belief-state reduction. Finite hidden states do not make the
space of probability beliefs a finite set.

\subsection{Find an executable plan satisfying task constraints}
\label{sec:p14-finder-planning}
\textit{Input and question.} An initial state, action operators, task network
and permitted decomposition methods: is there a finite permitted
decomposition that executes successfully?
\textit{Use and output.}
\hyperref[sec:p17-hierarchical-task-network-reasoning]{Hierarchical task
networks} specify decomposition, ordering, binding and intermediate-state
conditions. \cref{thm:p17-htn-finite-derivation} identifies the finite
derivation and execution witness.
\textit{Conditions.} A plan must satisfy the task-network constraints as well
as execute its actions. The general existence problem is semidecidable and
can be undecidable even under the restrictions in
\cref{thm:p17-htn-semidecision,thm:p17-htn-undecidable}.
For guaranteed finite expansion, the decreasing-level condition is
\cref{thm:p17-htn-levels}; no failed bounded expansion proves general
nonexistence.

\subsection{Assess a strategic profile}
\label{sec:p14-finder-games}
\textit{Input and question.} A finite normal-form game and candidate mixed
strategies: can any one player improve while the others stay fixed?
\textit{Use and output.}
\hyperref[def:p17-nash-equilibrium]{Nash equilibrium} asks exactly this
question. \cref{thm:p17-best-response-geometry} reduces its verification to
the stated pure-action payoff inequalities.
\textit{Conditions.} The expected-payoff calculation uses cardinal utilities
and independent mixed strategies. Equilibrium is a stability condition against
unilateral deviations, not a claim of collective optimality, uniqueness,
fairness or convergence of a learning process. The separate existence
constructions and their fixed-point foundations are in
\cref{sec:p17-finite-strategic-games}
\cite{Nash1950EquilibriumPoints,Nash1951NonCooperative}.

For an extensive game, \cref{sec:p17-rationality-9-1} distinguishes
contingent plans from behavior at information sets and preserves the
perfect-recall hypothesis of its realization/existence bridge. Finiteness
alone does not supply the corresponding behavioral result for an
absentminded game. For a population model,
\cref{sec:p17-rat93-evolutionary-games} separates invasion resistance, stationarity,
Lyapunov stability, attraction and mean payoff. Neither ESS nor a large
basin by itself establishes welfare or an ethical norm.

\subsection{Choose a collective procedure at its actual standard}
\label{sec:p14-finder-collective-standards}
\textit{Input and question.} Members' attitudes, a group procedure and its
output: is the criterion successful coordination, corporate consistency,
truth-directed reliability or justified communicative agreement?
\textit{Use and output.} \Cref{def:p17-rat102-diagnostic} separates
participant cooperation from corporate consistency;
\cref{def:p18-r101-selection} selects an epistemic procedure;
\cref{def:p18-r103-assessment} distinguishes recognition, uptake, truth
and institutional standing. \Cref{thm:p17-rat102-policy} supplies a
concrete feasible-individual/infeasible-majority policy contrast.
\textit{Conditions.} Shared agreement does not entail truth. An ordinal
cycle alone does not establish acceptance of priced trades. Jury reliability
requires the actual competence and dependence hypotheses of
\cref{thm:p18-r101-jury}; individual postdiscussion answers are not
independent group trials. Common belief is not factive common knowledge.

\subsection{Pool probabilities rather than binary judgments}
\label{sec:p14-finder-opinion-pooling}
\textit{Input and question.} Several probability assignments, pooling
weights and common evidence: does the specified aggregate remain coherent,
and how does it change under conditioning or partition refinement?
\textit{Use and output.} \Cref{def:p18-r101-pools,thm:p18-r101-pooling}
give finite linear/geometric constructions and their actual commutation
and refinement conditions. Their output is a probability assignment,
not a majority judgment set or preference ordering.
\textit{Conditions.} Check normalization, zeros, weights and the common
conditioning event. Coherence alone supplies neither independence of
the reports nor empirical accuracy. Identical pooled values do not identify
a deliberation mechanism \cite{DietrichSpiekermann2021Social}.

\subsection{Aggregate judgments on logically connected propositions}
\label{sec:p14-finder-aggregation}
\textit{Input and question.} An agenda, individual judgment sets and a
collective rule: are the resulting accepted propositions jointly coherent?
\textit{Use and output.}
\hyperref[sec:p17-judgment-aggregation]{Judgment aggregation} gives exact
agenda, closure and consistency conditions, the majority paradox, and
source-specific impossibility and restricted positive results.
\textit{Conditions.} Binary assent is different from a probability judgment
or preference ordering. Majority counts on each proposition separately need
not respect their logical dependencies; the concrete profile is
\cref{def:p17-ja-doctrinal-profile}. Check the selected theorem's electorate,
domain and rule axioms before applying an impossibility claim. Preference
aggregation has its own entry at \cref{sec:p17-social-choice}
\cite{ListPettit2001JudgmentPreprint}.

\subsection{Compute acceptability in an attack framework}
\label{sec:p14-finder-acceptability}
\textit{Input and question.} A finite list of arguments and attacks: which
arguments are acceptable under a specified extension semantics?
\textit{Use and output.}
\hyperref[def:p18-grounded-procedure]{Grounded iteration} returns its unique
extension; \cref{def:p18-extension-semantics} defines the separate complete,
preferred and stable alternatives.
\textit{Conditions.} Name the semantics and, for multiple extensions, whether
membership is required in some or every extension. Stable extensions may not
exist. The guarantees concern the supplied attack graph; they neither
establish premise truth nor justify the graph's attacks.
\cref{sec:p18-finite-argumentation} gives the constructions and separating
examples \cite[secs.~2.1--2.2, pp.~325--330]{Dung1995}.

\subsection{Decide what an unanswered challenge establishes}
\label{sec:p14-finder-challenge}
\textit{Input and question.} An argument, a critical question and the available
reply or evidence: has support been narrowed, suspended, defeated or rebutted?
\textit{Use and output.}
\cref{def:p18-scheme-assessment-terms,def:p18-scheme-reconstruction-method}
distinguish a request for support, an undercutting consideration and a reason
for the contrary conclusion.
\textit{Conditions.} Use the questions belonging to the selected expert,
witness, analogy, causal or practical scheme. Merely asking whether an
exception exists does not establish an exception; an unanswered question does
not automatically prove the opposite claim. For the different formal question
of whether answers are covered but still unsettled, use
\cref{thm:p18-qev-semantic-profile}.

\subsection{Interpret an empirical study of reasoning}
\label{sec:p14-finder-empirical}
\textit{Input and question.} A study report, task, sample and responses:
what empirical claim about reasoning do these records support?
\textit{Use and output.}
\hyperref[def:p19-evidence-record]{The empirical evidence record} connects
population, task version, scoring, design, estimate and uncertainty to a
bounded conclusion.
\textit{Conditions.} Separate the observed score from a proposed cognitive
process and the normative benchmark. Random recruitment and random assignment
address different questions. Repeated responses from one person are not
automatically independent participants. Read
\cref{sec:p19-evidence-methods} for the selected studies and interpretive
limits; this procedure is not itself an exhaustive literature review or an
accuracy certificate \cite[pp.~131--132]{WassersteinLazar2016ASA}.

\subsection{Interpret a developmental difference}
\label{sec:p14-finder-development}
\textit{Input and question.} Age groups or a child's responses under
specified tasks: does the observation concern inferential performance,
explicit understanding of validity, or confidence in an answer?
\textit{Use and output.} \Cref{def:p19-r151-explicit-implicit,
def:p19-r151-lookup} separate these constructs and preserve each
study's recruitment, language, instructions, scoring and denominators.
\textit{Conditions.} A familiar-task success does not establish general
transfer; an arbitrary-task failure does not prove absence of competence.
Cross-sectional age differences are not a longitudinal universal sequence.
Repeated rule-pattern records must not be counted as independent children
\cite{Markovits2021DevelopmentRationality}.

\subsection{Assess analogy and transfer to a new task}
\label{sec:p14-finder-transfer}
\textit{Input and question.} A source problem, proposed target problem and
mapping: which relations carry over, and is there evidence that people notice
and use them?
\textit{Use and output.}
\cref{sec:p12-heuristic-transformations-and-backward-reasoning} handles the
solution-preservation obligations of a mathematical transformation.
\cref{sec:p19-analogy-and-problem-transfer} separately treats the selected
analogy and transfer experiments and structural account.
\textit{Conditions.} Check mapped relations and target constraints; a useful
candidate still needs target verification. A formally available mapping does
not establish that participants retrieve it without a hint, or that a result
transfers across instructions, samples and domains
\cite{GickHolyoak1980AnalogicalProblemSolving,Gentner1983StructureMapping}.

\subsection{Evaluate a proposed reasoning intervention}
\label{sec:p14-finder-reasoning-intervention}
\textit{Input and question.} A training activity, comparison condition
and outcome: did it improve practiced performance, near or far transfer,
delayed retention, or decisions outside the tests?
\textit{Use and output.} \Cref{def:p19-r153-transfer,
def:p19-r153-intervention-method} specify those distinct achievements.
The original game/video and forecasting records retain their different
controls, selection mechanisms and follow-up populations.
\textit{Conditions.} Repeating an identical quiz immediately does not
test durable general transfer. A package comparison need not isolate
feedback, practice and duration, and selecting previous high performers
does not identify the causal effect of the selection procedure.
Measured accuracy under one benchmark does not by itself justify a
universal prescription \cite{DeOliveiraNisbett2021Improve,
MorewedgeEtAl2015Debiasing,MellersEtAl2014ForecastTraining}.

For comparisons after a task has been identified, use
\cref{sec:p14-decidability-and-complexity,
sec:p14-strength-consistency-and-independence,
sec:p14-expressiveness-and-representation}: their decision, strength,
expressiveness and correspondence tables retain the distinct questions
summarized here. The
selected method's own input conditions and output interpretation determine
what the resulting proof, model, estimate or recommendation establishes.
